% One A4 to hold at the lectern. Dense on purpose: the opening line of every slide is
% here verbatim, so if the memorized version deserts you mid-talk there is a sentence
% to grab rather than a keyword to reconstruct from.
%
% The full script is in the \note{} blocks of slides.tex -> slides-notes.pdf.
\documentclass[10pt,a4paper]{article}

\usepackage[T1]{fontenc}
\usepackage[utf8]{inputenc}
\usepackage[english]{babel}
\usepackage[margin=12mm,top=9mm,bottom=9mm]{geometry}
\usepackage{xcolor}
\usepackage{amsmath,amssymb}
\usepackage{array}

\definecolor{ctublue}{RGB}{0,101,189}
\definecolor{accent}{RGB}{200,16,46}

\renewcommand{\familydefault}{\sfdefault}
\setlength{\parindent}{0pt}
\setlength{\tabcolsep}{1.5mm}
\pagestyle{empty}

% opening line, said as written
\newcommand{\say}[1]{\textit{\color{black!75}``#1''}}
% a beat
\newcommand{\bt}{\newline\raisebox{0.15ex}{\tiny$\blacksquare$}~}
\newcommand{\cue}[3]{%
  \textbf{\color{ctublue}#1} & \textbf{#2} & #3 \\[1.5ex]}

\begin{document}

{\large\bfseries Acoustic Detection of Impulsive Sounds}\hfill
{\small 10 min -- target \textbf{9:30} -- \textbf{4 clicks}}

\vspace{0.5mm}\hrule\vspace{2mm}

{\small
\begin{tabular}{@{}>{\raggedright\arraybackslash}p{10mm} >{\raggedright\arraybackslash}p{19mm} >{\raggedright\arraybackslash}p{147mm}@{}}

\cue{0:00}{Title}{%
  \say{Good morning. I'm Martin Maxa from the Department of Measurement, here at this
  faculty. This is joint work with Jakub Svatoš from the same department.}
  \bt today: detecting, classifying, localizing impulsive sounds -- in our case, gunshots
  \bt \textbf{main part: gunshot classification} -- cepstral features $+$ two
  classifiers; localization briefly}

\cue{0:45}{Motivation}{%
  \say{This work is part of a larger project -- an acoustic sensor system for protecting
  soft targets.}
  \bt good complement to cameras: \textbf{omnidirectional}, \textbf{no line of sight}
  \bt \textbf{[CUT]} full system: detect, classify (to calibre), localize -- outdoors,
  noise, a lot of reflections
  \bt \textbf{\color{accent}[CLICK]} \textbf{three questions:} (1) which cepstral
  features best separate gunshots from similar impulsive sounds, (2) NN or SVM,
  (3) frame length vs.\ results}

\cue{1:50}{Signature}{%
  \say{Let's start with the signal of a gunshot. A gunshot has one or two components,
  depending on the projectile speed.}
  \bt \textbf{subsonic $\to$ muzzle blast only; supersonic $\to$ also a shock wave}
  \say{this recording is a supersonic shot, so we can see both}
  \bt \textbf{shock wave} $\approx$2.5\,ms -- \textbf{[CUT]} along the bullet's path, typical
  N-shape
  \bt \textbf{\color{accent}[CLICK]} \textbf{muzzle blast} $\approx$45\,ms later -- gases
  leaving the barrel, \emph{every gunshot has it}
  \bt between and after $=$ \textbf{reflections} -- \say{this matters for the frame length,
  and I'll come back to it in the results}}

\cue{3:00}{Single-domain}{%
  \say{This also shows why one domain alone is not enough.}
  \bt spectrum $\to$ tells you more about the environment than about the gunshot
  \bt waveform $\to$ changes as the sound travels and hits obstacles
  \bt $\Rightarrow$ features combining both domains -- cepstral features do exactly that}

\cue{3:30}{System}{%
  \say{Now to the system itself. It's a set of independent sensor units connected to a
  remote server.}
  \bt unit monitors continuously $\to$ on an impulsive event, sends the signal
  \bt server does two things -- \textbf{classifies} the event, \textbf{localizes} the source
  \bt more units: wider coverage $+$ position of the source $|$ tested with several
  calibres and types of ammunition}

\cue{4:20}{Localize}{%
  \say{Let me briefly explain the localization first. Each sensor unit has two MEMS
  microphones, about twenty centimetres apart, sampled at 44.1 kilohertz.}
  \bt TDoA between the microphones $\to$ $\Theta = \arcsin\!\big(c\,\Delta n / (d f_s)\big)$
  $\to$ direction to the source
  \bt \textbf{\color{accent}[CLICK]} one unit $=$ one direction; two units $\to$ where they
  cross $=$ position
  \bt \textbf{70 events}, 9 reference angles -- cross-corr $+$ parabolic interpolation,
  \textbf{MAE $0.77^\circ$}}

\cue{5:20}{Features}{%
  \say{Now to the main part -- classification. We compared four cepstral feature
  extractors. They differ in the filter bank they use.}
  \bt MFCC mel, dense low (from speech; works for gunshots too) $|$ LFCC linear $|$ IMFCC
  inverted mel, dense high $|$ \textbf{\color{accent}GTCC gammatone}, model of human hearing,
  more noise-robust $|$ (grid $4\times3\times2$ is on the slide)
  \bt \textbf{\color{accent}NEXT} (5:55) \say{And this is how the features are computed.}
  same chain: framing, FFT, \textbf{filter bank}, log, DCT $|$ \textbf{filter bank $=$ the only
  difference} $|$ coefficients $\to$ \textbf{classifier: NN and SVM}, frames 15/30/50\,ms}

\cue{6:25}{\color{accent}TABLE $\to$ CHART}{%
  \say{This is the full table of results for the nine-millimetre class.}
  \bt GTCC rows highlighted $|$ \emph{don't read the numbers} $\to$ \textbf{\color{accent}NEXT} (6:40)
  \say{Here are the same results, just the Matthews correlation coefficient.}
  \bt \textbf{MCC:} whole confusion matrix, less affected by class imbalance
  \bt \textbf{red $=$ GTCC on top in both panels, every frame length} $|$ \textbf{NN does
  better than SVM, gap grows with longer frames}
  \bt best \textbf{GTCC $+$ NN, 50\,ms $\to$ 99.15 \,/\, 96.77 \,/\, 0.98} $\Rightarrow$
  \textbf{117 of 118} correct (31 of them 9\,mm; 30 of 31 recalled)}

\cue{7:35}{Discussion}{%
  \say{Three things to take from this.}
  \bt MFCC, LFCC, IMFCC similar -- filter spacing makes little difference
  \bt \textbf{GTCC clearly better} -- \textbf{[CUT]} hearing model, noise-robust:
  \emph{likely} reason, not tested directly
  \bt \textbf{NN slightly better than SVM}, especially longer frames -- \emph{general trend,
  not every single result}
  \bt \textbf{\color{accent}[CLICK]} longer frames improve accuracy, \emph{but} at 50\,ms the
  frame already contains more and more reflections -- \say{so even longer frames wouldn't
  necessarily help}}

\cue{8:45}{Conclusion}{%
  \say{To sum up.}
  \bt distributed sensor system $\to$ detects, classifies, localizes impulsive sounds
  \bt cepstral features work well with both classifiers
  \bt GTCC best of the four; NN slightly better than SVM, especially with longer frames
  \bt all tested on real firearms, several calibres and types of ammunition}

\cue{9:30}{}{\say{Thank you.}}

\end{tabular}
}

\vspace{-1mm}\hrule\vspace{1.5mm}

{\footnotesize
\textbf{\color{accent}If asked, the numbers.}
Corpus 371 train \,/\, 118 test, five classes (false alarms 86/21, \textbf{9\,mm 70/31},
5.56 NATO SD 100/31, 7.62 Tokarev 55/17, .22 60/18) $|$ 26 features per extractor,
$\approx$75/25 split $|$ NN: MATLAB Neural Fitting, Levenberg--Marquardt, 370 iterations $|$
bearing: cross-corr $0.84^\circ$, $+$parabolic $\mathbf{0.77^\circ}$, PBDE $3.24^\circ$,
PBDE$+$linreg $3.91^\circ$ $|$ detection stage $=$ median-filter trigger, RTSI 2019, down to
SNR $\approx$5\,dB $|$ appendix $=$ last 4 PDF pages: \textbf{A1}~setup \textbf{A2}~bearing
\textbf{A3}~other~cepstral \textbf{A4}~metrics (RLC $=$ recall).
\par\vspace{1mm}
\textbf{\color{accent}Running late?} Three passages are marked \textbf{[CUT IF LATE]} in the
notes -- \textbf{Motivation} (the three functions), \textbf{Signature} (Mach cone),
\textbf{Discussion} (probable reason GTCC wins). All three $\approx$35\,s; dropping them
breaks nothing. \textbf{Checkpoint:} start the table at \textbf{6:25}; past 6:55, cut all three.
\par\vspace{1mm}
\textbf{\color{accent}Do not guess:} 2D position error in metres, SNR limits of \emph{this}
classifier, firing distances in this dataset, provenance of the waveform on slide 3.
\say{We did not characterize that in this study -- Jakub Svatoš ran the campaign, I can
put you in touch.}}

\end{document}
